\documentclass[10pt,twocolumn]{article}
\usepackage[margin=0.75in]{geometry}
\usepackage{amsmath}
\usepackage{booktabs}
\usepackage{hyperref}
\usepackage{graphicx}
\usepackage{listings}
\usepackage{xcolor}
\usepackage{microtype}
\usepackage{parskip}
\usepackage{titlesec}
\usepackage{abstract}

\titleformat{\section}{\normalfont\large\bfseries}{\thesection}{1em}{}
\titleformat{\subsection}{\normalfont\normalsize\bfseries}{\thesubsection}{1em}{}

\lstset{
  basicstyle=\ttfamily\footnotesize,
  breaklines=true,
  frame=single,
  backgroundcolor=\color{gray!10},
}

\hypersetup{
  colorlinks=true,
  linkcolor=black,
  citecolor=black,
  urlcolor=blue,
}

\title{\textbf{ArchProbe: Measuring Cross-Architecture Stability of\\
Syscall Behavioral Signatures in Containerized Environments}}

\author{
  Sherbekjon Rustamov \\
  Independent Security Researcher \\
  \texttt{rustamovsherbekjon@gmail.com}
}

\date{}

\begin{document}

\maketitle

\begin{abstract}
Syscall-based behavioral detection systems are trained and deployed
predominantly on x86 infrastructure, yet modern cloud and edge deployments
span both x86 and ARM architectures. This paper asks whether a classifier
trained on x86 syscall traces can generalize to ARM without retraining.
We present \textbf{ArchProbe}, a controlled measurement study that collects
\texttt{strace} behavioral traces from 20 containerized workloads (10 benign,
10 malicious-pattern) on both \texttt{x86\_64} and \texttt{aarch64} platforms,
then evaluates four cross-architecture transfer experiments using a Random
Forest classifier. Our results show that raw syscall frequency features
collapse to chance-level performance when applied cross-architecture
(50\% accuracy, F1\,=\,0.00), while an abstract category normalization
layer---mapping architecture-specific syscall names to semantic
groups---recovers 65\% accuracy (F1\,=\,0.53). These findings expose a
silent assumption in syscall-based intrusion detection: architectural
portability is not free, and normalization is a necessary but insufficient
remedy at small data scales.
\end{abstract}

% -----------------------------------------------------------------------
\section{Introduction}
% -----------------------------------------------------------------------

Syscall-based intrusion detection has a long history in systems security.
Tools such as Falco, Sysdig, and academic systems like Optimus~\cite{optimus}
operate by observing the stream of system calls a process makes and classifying
whether the pattern is benign or malicious. The implicit assumption in every
deployed system is that the model trained on one machine generalizes to others
running the same workload.

This assumption silently breaks across CPU architectures. An \texttt{x86\_64}
Linux binary and its \texttt{aarch64} equivalent may use different system calls
to accomplish the same task: \texttt{mmap} vs.\ \texttt{mmap2}, \texttt{stat}
vs.\ \texttt{newfstatat}, or \texttt{clone} vs.\ \texttt{clone3}. ARM's
expanded syscall table and different calling conventions mean that the raw
string identifiers of syscalls differ even for identical logical operations. A
classifier trained to recognize \texttt{privesc\_enum} patterns on x86 by the
frequency of \texttt{openat}, \texttt{getdents64}, and \texttt{readlinkat}
may encounter entirely different call sequences for the same attack on ARM.

With ARM server adoption accelerating---AWS Graviton instances, Ampere Altra
deployments, and Apple Silicon development machines---this assumption is
becoming load-bearing infrastructure debt. ArchProbe is a measurement study
designed to quantify the failure mode and test a normalization remedy.

\textbf{Research questions:}
\begin{itemize}
  \item \textbf{RQ1:} Do containerized workloads produce
    architecture-invariant syscall behavioral signatures when represented
    as raw call frequencies?
  \item \textbf{RQ2:} Can mapping raw syscall names to abstract semantic
    categories recover cross-architecture classification accuracy?
\end{itemize}

% -----------------------------------------------------------------------
\section{Related Work}
% -----------------------------------------------------------------------

\textbf{Syscall-based anomaly detection.}
The use of system call sequences for intrusion detection dates to Forrest et
al.~\cite{forrest1996}, who proposed normal behavior profiles from syscall
sequences. Subsequent work explored bag-of-syscalls, n-gram models, and neural
sequence models~\cite{abed2015}. These techniques underpin modern container
security tools.

\textbf{Optimus.}
Optimus~\cite{optimus} reduces container attack surface by statically filtering
syscalls a container is permitted to make. It derives per-architecture policies
independently. Our work asks what happens when a \emph{learned} policy---rather
than a static allowlist---must transfer across architectures.

\textbf{Cross-architecture analysis.}
Prior work on architecture-independent malware analysis has focused on binary
similarity (instruction embedding, control-flow graph matching). To our
knowledge, no prior work has directly measured syscall-level behavioral
portability for \emph{containerized} workloads and quantified the resulting
accuracy degradation.

% -----------------------------------------------------------------------
\section{Methodology}
% -----------------------------------------------------------------------

\subsection{Testbed}

All traces are collected inside Docker containers using
\texttt{strace -f -e trace=all}. The host system is FothOS (an Arch Linux
security distribution). Docker Buildx enables building and running
\texttt{linux/arm64} images on the same \texttt{x86\_64} host via QEMU
userspace emulation, ensuring that ARM traces reflect \texttt{aarch64} binary
behavior under the Linux aarch64 ABI.

\subsection{Workloads}

We define 20 workloads: 10 benign and 10 malicious-pattern. All malicious
workloads are fully simulated inside isolated containers with no real external
targets.

\begin{table}[h]
\centering
\small
\begin{tabular}{ll}
\toprule
\textbf{Category} & \textbf{Workloads} \\
\midrule
Benign (10) & curl, DNS resolution, file I/O, \\
            & filesystem find, process fork, \\
            & git clone, proc.\ enum., Python, \\
            & SSH keygen, wget \\
\midrule
Malicious (10) & C2 beacon, cred.\ harvest, \\
               & data exfil, file staging, \\
               & log wipe, persistence, \\
               & privesc enum., net.\ recon., \\
               & proc.\ recon., reverse shell sim. \\
\bottomrule
\end{tabular}
\caption{Workload categories}
\end{table}

\textbf{Data quality note.} Five \texttt{x86\_64} benign traces (dns, find,
python, ssh\_keygen, wget) produced empty output, likely due to QEMU
initialization timing during \texttt{strace} attachment on short-lived
processes. These are excluded, leaving 15 x86 records (5 benign, 10 malicious)
and 20 ARM records (10 benign, 10 malicious).

\subsection{Feature Representations}

\textbf{Frequency vector.} A bag-of-syscalls: \texttt{\{syscall: count\}},
normalized to relative frequencies. Lines matching
\texttt{syscall\_0x[0-9a-f]+} are excluded as unresolved kernel pointer noise.

\textbf{Normalized categories.} Each syscall name is mapped to one of eight
semantic groups: \texttt{file\_io}, \texttt{network}, \texttt{process},
\texttt{memory}, \texttt{privilege}, \texttt{ipc}, \texttt{timing},
\texttt{sync}. Unknown syscalls map to \texttt{other}. The mapping is
architecture-agnostic: both \texttt{mmap} and \texttt{mmap2} map to
\texttt{memory}; both \texttt{open} and \texttt{openat} map to
\texttt{file\_io}.

\subsection{Classifier}

We use a Random Forest classifier (100 trees, \texttt{class\_weight="balanced"},
\texttt{random\_state=42}) from scikit-learn. Balanced class weights compensate
for the 5:10 benign:malicious imbalance in x86 training data.

\subsection{Experiments}

\begin{table}[h]
\centering
\small
\begin{tabular}{clll}
\toprule
\textbf{ID} & \textbf{Train} & \textbf{Test} & \textbf{Feature} \\
\midrule
A & x86 (60\%, stratified) & x86 (40\%) & frequency \\
B & x86 (all 15) & ARM (all 20) & frequency \\
C & x86 + ARM first half & ARM second half & frequency \\
D & x86 (all 15) & ARM (all 20) & normalized \\
\bottomrule
\end{tabular}
\caption{Experiment configurations. B and D are the primary cross-architecture
transfer tests.}
\end{table}

% -----------------------------------------------------------------------
\section{Results}
% -----------------------------------------------------------------------

\begin{table}[h]
\centering
\small
\begin{tabular}{lrr}
\toprule
\textbf{Experiment} & \textbf{Accuracy} & \textbf{F1} \\
\midrule
A: x86$\to$x86 (baseline)          & 100.0\% & 1.0000 \\
B: x86$\to$ARM (raw features)      &  50.0\% & 0.0000 \\
C: x86+ARM$\to$ARM (mixed)         &   0.0\% & 0.0000 \\
D: x86$\to$ARM (normalized)        &  65.0\% & 0.5333 \\
\bottomrule
\end{tabular}
\caption{Classification results across four experiments.}
\end{table}

\textbf{Experiment B} confusion matrix:

\begin{center}
\small
\begin{tabular}{lrr}
\toprule
 & Pred.\ Benign & Pred.\ Malicious \\
\midrule
True Benign    & 10 & 0 \\
True Malicious & 10 & 0 \\
\bottomrule
\end{tabular}
\end{center}

The classifier predicts every ARM sample as benign. Top features
(\texttt{poll}, \texttt{read}, \texttt{write}, \texttt{open}) are common x86
syscalls that are absent or infrequent in ARM traces under QEMU, leaving no
discriminating signal.

\textbf{Experiment D} confusion matrix:

\begin{center}
\small
\begin{tabular}{lrr}
\toprule
 & Pred.\ Benign & Pred.\ Malicious \\
\midrule
True Benign    & 9 & 1 \\
True Malicious & 6 & 4 \\
\bottomrule
\end{tabular}
\end{center}

Four malicious ARM samples are correctly identified. Top features:
\texttt{ipc} (0.29), \texttt{file\_io} (0.23), \texttt{other} (0.16),
\texttt{memory} (0.12), \texttt{timing} (0.08). These are architecture-agnostic
categories, which is why any transfer occurs at all.

\textbf{Experiment C} (0\% accuracy) suffers from a structural data artifact:
positional splitting of ARM records by filename sort order places all ARM benign
samples in training and all malicious in test. This result is excluded from the
primary findings.

% -----------------------------------------------------------------------
\section{Discussion}
% -----------------------------------------------------------------------

\subsection{RQ1: Raw Frequency Features}

Raw syscall frequency features do not transfer across architectures.
Experiment B shows complete collapse: the classifier defaults to predicting
the majority class across all 20 ARM samples.

This finding has a direct operational implication: a model trained on x86
telemetry and applied to ARM workloads without adaptation is not merely less
accurate---it is effectively disabled, operating at chance level. Deployed
detection systems targeting multi-architecture fleets face a silent blind spot
on every ARM node.

\subsection{RQ2: Normalized Category Features}

Normalization to abstract categories partially recovers cross-architecture
accuracy (65\%, F1\,=\,0.53). The \texttt{ipc} and \texttt{file\_io}
categories are the strongest discriminating features. This makes intuitive
sense: malicious workloads performing network reconnaissance or data
exfiltration use IPC-class and file I/O calls heavily regardless of
architecture, because workload \emph{intent} determines category distributions
more than architecture does.

However, 65\% accuracy is not production-viable. Six of ten malicious ARM
workloads are missed. Normalization is necessary but not sufficient at this
data scale.

\subsection{Limitations}

\begin{enumerate}
  \item \textbf{Small dataset.} 15 x86 training samples is insufficient for
    robust generalization. Results are directional, not definitive.
  \item \textbf{QEMU emulation.} ARM traces were collected under QEMU
    userspace emulation, not native \texttt{aarch64} hardware. QEMU's
    translation layer may introduce artifacts absent on real ARM.
  \item \textbf{Simulated malice.} All malicious-pattern workloads are
    synthetic. Real malware behavioral signatures may differ.
  \item \textbf{Single classifier.} Only Random Forest was evaluated.
    Sequence-aware models (HMM, LSTM) may behave differently.
\end{enumerate}

\subsection{Practical Recommendations}

\begin{enumerate}
  \item \textbf{Do not assume portability.} Publish per-architecture model
    versions or include architecture as an explicit feature.
  \item \textbf{Normalize before training.} Abstract syscall categories are
    a lightweight, interpretable normalization layer that improves
    cross-architecture transfer without retraining.
  \item \textbf{Collect native ARM data.} QEMU emulation is not a substitute
    for native traces; models trained on native x86 + native ARM data should
    outperform the mixed QEMU results reported here.
\end{enumerate}

% -----------------------------------------------------------------------
\section{Conclusion}
% -----------------------------------------------------------------------

ArchProbe demonstrates that syscall behavioral signatures are not
architecture-invariant. A Random Forest classifier trained on x86 container
traces achieves 100\% accuracy on held-out x86 data but collapses to chance
(50\%, F1\,=\,0.00) when applied directly to ARM traces. Mapping syscall names
to abstract semantic categories partially recovers transfer capability (65\%,
F1\,=\,0.53), confirming that workload intent is preserved across architectures
even when raw call names diverge.

This work motivates further research into architecture-aware behavioral
detection: larger multi-architecture datasets, native ARM collection, and
sequence-aware normalization. The ArchProbe dataset and code are publicly
available at \url{https://github.com/SherbekjonDev/ArchProbe}.

% -----------------------------------------------------------------------
\bibliographystyle{plain}
\begin{thebibliography}{9}

\bibitem{optimus}
Z.~Wan et al.,
``Sysfilter: Automated System Call Filtering for Commodity Software,''
\emph{RAID}, 2020.

\bibitem{forrest1996}
S.~Forrest, S.~A.~Hofmeyr, A.~Somayaji, and T.~A.~Longstaff,
``A sense of self for Unix processes,''
\emph{IEEE Symposium on Security and Privacy}, 1996.

\bibitem{abed2015}
A.~S.~Abed, C.~Clancy, and D.~S.~Levy,
``Applying bag of system calls for anomalous behavior detection of applications
in Linux containers,''
\emph{IEEE Global Humanitarian Technology Conference}, 2015.

\bibitem{falco}
Falco Project, ``Cloud-native runtime security,''
\url{https://falco.org}, 2024.

\end{thebibliography}

\end{document}
